\documentclass[11pt,a4paper]{article}
% Preambulo comun de los reportes de Tecnologias Celulares.
% Cada main.tex hace: \documentclass ... % Preambulo comun de los reportes de Tecnologias Celulares.
% Cada main.tex hace: \documentclass ... % Preambulo comun de los reportes de Tecnologias Celulares.
% Cada main.tex hace: \documentclass ... \input{../comun/preambulo.tex}
\usepackage[utf8]{inputenc}
\usepackage[T1]{fontenc}
\usepackage[spanish,mexico]{babel}
\usepackage{graphicx}
\usepackage[margin=2.5cm]{geometry}
\usepackage{amsmath}
\usepackage{amssymb}
\usepackage{booktabs}
\usepackage{tabularx}
\usepackage{array}
\usepackage{float}
\usepackage{setspace}
\usepackage{parskip}
\usepackage{fancyhdr}
\usepackage{textcomp}
\usepackage{xcolor}
\usepackage{caption}
\usepackage{fvextra}
\usepackage[hidelinks]{hyperref}

\onehalfspacing
\setlength{\parskip}{6pt}
\sloppy
\setlength{\emergencystretch}{3em}
\setlength{\headheight}{40pt}
\captionsetup{font=small,labelfont=bf}
\setlength{\tabcolsep}{6pt}

\pagestyle{fancy}
\fancyhf{}
\fancyhead[L]{\includegraphics[width=1cm]{/home/yovick/repos/UAZ-IEI/Assets/UIE.png}}
\fancyhead[R]{\small Tecnologías Celulares}
\fancyfoot[C]{\thepage}
\renewcommand{\headrulewidth}{0.4pt}

\DefineVerbatimEnvironment{codigo}{Verbatim}{%
  breaklines=true,breakanywhere=true,fontsize=\small,frame=single,
  framesep=3pt,rulecolor=\color{gray!50}}

% Figura centrada con pie: \figura[ancho]{archivo}{leyenda}
\newcommand{\figura}[3][0.85]{%
  \begin{center}
    \includegraphics[width=#1\linewidth]{#2}\\[2pt]
    {\small #3}
  \end{center}}

% Portada: \portada{numero}{titulo de la practica}
\newcommand{\portada}[2]{%
\begin{titlepage}
  \thispagestyle{empty}
  \begin{center}
    \vspace*{1cm}
    \includegraphics[width=4.5cm]{/home/yovick/repos/UAZ-IEI/Assets/uaz.png}
    \vspace{0.8cm}

    {\large Universidad Autónoma de Zacatecas}\\[4pt]
    {\large Unidad Académica de Ingeniería Eléctrica}\\[4pt]
    {\normalsize Ingeniería Electrónica Industrial con especialidad en Telecomunicaciones}

    \vspace{1.8cm}
    {\LARGE\bfseries Práctica #1}\\[10pt]
    {\Large\bfseries #2}
  \end{center}

  \vspace{2.2cm}
  \begin{center}
    \begin{tabular}{rl}
      \textbf{Materia:} & Tecnologías Celulares\\
      \textbf{Docente:} & Dr. Víktor Iván Rodríguez Abdalá\\
      \textbf{Alumno:}  & Arnold Yovick Ríos Zaldívar\\
      \textbf{Fecha:}   & Noviembre de 2026\\
    \end{tabular}
  \end{center}
  \vfill
\end{titlepage}}

\usepackage[utf8]{inputenc}
\usepackage[T1]{fontenc}
\usepackage[spanish,mexico]{babel}
\usepackage{graphicx}
\usepackage[margin=2.5cm]{geometry}
\usepackage{amsmath}
\usepackage{amssymb}
\usepackage{booktabs}
\usepackage{tabularx}
\usepackage{array}
\usepackage{float}
\usepackage{setspace}
\usepackage{parskip}
\usepackage{fancyhdr}
\usepackage{textcomp}
\usepackage{xcolor}
\usepackage{caption}
\usepackage{fvextra}
\usepackage[hidelinks]{hyperref}

\onehalfspacing
\setlength{\parskip}{6pt}
\sloppy
\setlength{\emergencystretch}{3em}
\setlength{\headheight}{40pt}
\captionsetup{font=small,labelfont=bf}
\setlength{\tabcolsep}{6pt}

\pagestyle{fancy}
\fancyhf{}
\fancyhead[L]{\includegraphics[width=1cm]{/home/yovick/repos/UAZ-IEI/Assets/UIE.png}}
\fancyhead[R]{\small Tecnologías Celulares}
\fancyfoot[C]{\thepage}
\renewcommand{\headrulewidth}{0.4pt}

\DefineVerbatimEnvironment{codigo}{Verbatim}{%
  breaklines=true,breakanywhere=true,fontsize=\small,frame=single,
  framesep=3pt,rulecolor=\color{gray!50}}

% Figura centrada con pie: \figura[ancho]{archivo}{leyenda}
\newcommand{\figura}[3][0.85]{%
  \begin{center}
    \includegraphics[width=#1\linewidth]{#2}\\[2pt]
    {\small #3}
  \end{center}}

% Portada: \portada{numero}{titulo de la practica}
\newcommand{\portada}[2]{%
\begin{titlepage}
  \thispagestyle{empty}
  \begin{center}
    \vspace*{1cm}
    \includegraphics[width=4.5cm]{/home/yovick/repos/UAZ-IEI/Assets/uaz.png}
    \vspace{0.8cm}

    {\large Universidad Autónoma de Zacatecas}\\[4pt]
    {\large Unidad Académica de Ingeniería Eléctrica}\\[4pt]
    {\normalsize Ingeniería Electrónica Industrial con especialidad en Telecomunicaciones}

    \vspace{1.8cm}
    {\LARGE\bfseries Práctica #1}\\[10pt]
    {\Large\bfseries #2}
  \end{center}

  \vspace{2.2cm}
  \begin{center}
    \begin{tabular}{rl}
      \textbf{Materia:} & Tecnologías Celulares\\
      \textbf{Docente:} & Dr. Víktor Iván Rodríguez Abdalá\\
      \textbf{Alumno:}  & Arnold Yovick Ríos Zaldívar\\
      \textbf{Fecha:}   & Noviembre de 2026\\
    \end{tabular}
  \end{center}
  \vfill
\end{titlepage}}

\usepackage[utf8]{inputenc}
\usepackage[T1]{fontenc}
\usepackage[spanish,mexico]{babel}
\usepackage{graphicx}
\usepackage[margin=2.5cm]{geometry}
\usepackage{amsmath}
\usepackage{amssymb}
\usepackage{booktabs}
\usepackage{tabularx}
\usepackage{array}
\usepackage{float}
\usepackage{setspace}
\usepackage{parskip}
\usepackage{fancyhdr}
\usepackage{textcomp}
\usepackage{xcolor}
\usepackage{caption}
\usepackage{fvextra}
\usepackage[hidelinks]{hyperref}

\onehalfspacing
\setlength{\parskip}{6pt}
\sloppy
\setlength{\emergencystretch}{3em}
\setlength{\headheight}{40pt}
\captionsetup{font=small,labelfont=bf}
\setlength{\tabcolsep}{6pt}

\pagestyle{fancy}
\fancyhf{}
\fancyhead[L]{\includegraphics[width=1cm]{/home/yovick/repos/UAZ-IEI/Assets/UIE.png}}
\fancyhead[R]{\small Tecnologías Celulares}
\fancyfoot[C]{\thepage}
\renewcommand{\headrulewidth}{0.4pt}

\DefineVerbatimEnvironment{codigo}{Verbatim}{%
  breaklines=true,breakanywhere=true,fontsize=\small,frame=single,
  framesep=3pt,rulecolor=\color{gray!50}}

% Figura centrada con pie: \figura[ancho]{archivo}{leyenda}
\newcommand{\figura}[3][0.85]{%
  \begin{center}
    \includegraphics[width=#1\linewidth]{#2}\\[2pt]
    {\small #3}
  \end{center}}

% Portada: \portada{numero}{titulo de la practica}
\newcommand{\portada}[2]{%
\begin{titlepage}
  \thispagestyle{empty}
  \begin{center}
    \vspace*{1cm}
    \includegraphics[width=4.5cm]{/home/yovick/repos/UAZ-IEI/Assets/uaz.png}
    \vspace{0.8cm}

    {\large Universidad Autónoma de Zacatecas}\\[4pt]
    {\large Unidad Académica de Ingeniería Eléctrica}\\[4pt]
    {\normalsize Ingeniería Electrónica Industrial con especialidad en Telecomunicaciones}

    \vspace{1.8cm}
    {\LARGE\bfseries Práctica #1}\\[10pt]
    {\Large\bfseries #2}
  \end{center}

  \vspace{2.2cm}
  \begin{center}
    \begin{tabular}{rl}
      \textbf{Materia:} & Tecnologías Celulares\\
      \textbf{Docente:} & Dr. Víktor Iván Rodríguez Abdalá\\
      \textbf{Alumno:}  & Arnold Yovick Ríos Zaldívar\\
      \textbf{Fecha:}   & Noviembre de 2026\\
    \end{tabular}
  \end{center}
  \vfill
\end{titlepage}}


\begin{document}

\portada{12}{Proyecto integrador: caracterización experimental de una red LTE utilizando RTL-SDR y MATLAB}

\textbf{Script:} \texttt{practica12.py} (Python con NumPy/SciPy)

\section{Resumen}

Se caracterizó una red LTE en la banda de 850 MHz a partir de capturas IQ de 262\,144
muestras (17.1 ms) en tres ubicaciones: exterior del edificio, interior del laboratorio y
pasillo. Para cada una se estimaron la potencia recibida, el piso de ruido, la relación
señal a ruido, el ancho de banda ocupado y la capacidad aproximada. La potencia cayó de
$-$29.16 dBFS en el exterior a $-$47.29 dBFS en el laboratorio y $-$54.30 dBFS en el
pasillo, con SNR de 40.7, 21.9 y 15.9 dB y capacidades de 188.3, 102.7 y 74.2 Mbps. La
multitrayectoria del pasillo fragmentó el espectro de las portadoras, evidencia directa de
la dispersión del canal. Las capturas son sintéticas porque no había RTL-SDR conectado; el
procesamiento es el mismo que se aplicaría a una captura real.

\section{Introducción}

LTE utiliza OFDM en el enlace descendente y técnicas MIMO para aumentar la capacidad y la
eficiencia espectral. Caracterizar una red real implica medir su ocupación espectral, su
potencia y su relación señal a ruido en distintos puntos, y relacionar esas medidas con los
conceptos de propagación, dispersión y capacidad estudiados durante el curso.

Este proyecto integra las prácticas anteriores: la adquisición y el análisis espectral de
la práctica 1, el modelo OFDM de las prácticas 5 y 6, y la capacidad MIMO de la práctica 9.

\section{Marco teórico}

\subsection{Banda de 850 MHz y estructura LTE}

En México, la banda celular de 850 MHz utiliza el bloque de enlace descendente
869--894 MHz. Un canal LTE ocupa un ancho nominal de 1.4, 3, 5, 10, 15 o 20 MHz y se
construye con subportadoras de 15 kHz. Para un ancho de 10 MHz se usan 600 subportadoras
(9 MHz ocupados) y para 5 MHz, 300 subportadoras (4.5 MHz).

La señal LTE es multiportadora: su densidad espectral de potencia es aproximadamente plana
dentro del canal y cae abruptamente en los bordes, con bandas de guarda entre canales. Esa
forma es la firma espectral de OFDM y permite identificar las portadoras con una FFT.

\subsection{Propagación y capacidad}

La potencia recibida disminuye con la distancia y con los obstáculos. Las pérdidas de
penetración en interiores y la atenuación adicional de un pasillo explican las diferencias
entre ubicaciones. La multitrayectoria produce desvanecimientos selectivos en frecuencia:
ciertas subportadoras caen en un nulo y otras no, lo que fragmenta el espectro observado.

La capacidad aproximada de un canal AWGN de ancho $B$ es
\begin{equation}
  C = B\log_2\left(1+\mathrm{SNR}\right),
\end{equation}
donde la SNR se estima a partir de la PSD medida. Para una red MIMO con $N_t$ antenas
transmisoras y $N_r$ receptoras, la capacidad puede crecer hasta en un factor
$\min(N_t,N_r)$ respecto de SISO.

\section{Metodología}

\subsection{Adquisición}

\begin{tabularx}{\linewidth}{@{}*{2}{>{\raggedright\arraybackslash}X}@{}}
  \toprule
  Parámetro & Valor \\
  \midrule
  Frecuencia central & 881 MHz \\
  Frecuencia de muestreo & 15.36 MSPS (1024 FFT, 15 kHz por subportadora) \\
  Muestras & 262\,144 (17.1 ms) \\
  Portadoras & 10 MHz en 878 MHz y 5 MHz en 885.5 MHz \\
  Receptor & RTL-SDR (captura sintética; el equipo no estaba conectado) \\
  \bottomrule
\end{tabularx}

Las capturas se generaron con dos portadoras OFDM de 15 kHz por subportadora, un canal
multitrayectoria de tres ecos y ruido gaussiano. Se aplicó un filtro de banda con flancos
de coseno alzado para reproducir la emisión fuera de banda de un transmisor real. No se
incluyó el prefijo cíclico porque el análisis es espectral, no de demodulación.

\subsection{Procesamiento}

\begin{enumerate}
  \item PSD de Welch con 4096 muestras por segmento y 50\,\% de solapamiento.
  \item Piso de ruido estimado con el percentil 5 de la PSD (las portadoras ocupan la mayor
        parte del ancho de banda).
  \item Detección de regiones 8 dB por encima del piso; ancho de banda al 99\,\% de la
        potencia por integración de la PSD.
  \item Potencia total $P=\frac{1}{N}\sum|x[n]|^2$ en dBFS.
  \item SNR en banda a partir de la potencia de las regiones detectadas y el piso de ruido.
  \item Capacidad con $C=B\log_2(1+\mathrm{SNR})$.
\end{enumerate}

\section{Escenarios de medición}

\begin{tabularx}{\linewidth}{@{}l>{\raggedright\arraybackslash}Xcc@{}}
  \toprule
  Ubicación & Descripción & Potencia relativa & Ecos \\
  \midrule
  Exterior & Explanada frente al edificio, línea de vista a la estación base &
  $-$30 / $-$38 dBFS & 0, 1 y 3 muestras \\
  Laboratorio & Interior, paredes y equipo, sin línea de vista &
  $-$48 / $-$56 dBFS & 0, 4 y 9 muestras \\
  Pasillo & Corredor largo, reflexiones múltiples &
  $-$55 / $-$63 dBFS & 0, 7 y 15 muestras \\
  \bottomrule
\end{tabularx}

La potencia indicada corresponde a las dos portadoras (10 y 5 MHz) antes del canal. El
canal multitrayectoria se normaliza a energía unitaria y el ruido térmico se fija en
$-$72 dBFS en las tres ubicaciones.

\section{Resultados experimentales}

\subsection{Espectros}

\figura[0.98]{figs/espectros.png}{PSD de las tres ubicaciones. La línea punteada es el piso de ruido estimado y las sombras marcan las regiones detectadas.}

En el exterior se distinguen con claridad las dos portadoras, con una meseta de
9 MHz y otra de 4.5 MHz separadas por una guarda; la ligera inclinación de la meseta es
la respuesta en frecuencia del canal, no de la emisión. En el laboratorio la meseta aparece
ondulada por los desvanecimientos selectivos. En el pasillo los ecos producen nulos
profundos que fragmentan cada portadora en varias regiones.

\subsection{Espectrogramas}

\figura[0.98]{figs/espectrogramas.png}{Espectrogramas de las tres ubicaciones.}

Las tres capturas son estacionarias en la ventana de 17 ms: no hay saltos de frecuencia ni
ráfagas. La diferencia entre ubicaciones está en el nivel y en la forma del espectro, no en
la estabilidad temporal.

\subsection{Mediciones}

\begin{tabular}{@{}lcccccc@{}}
  \toprule
  Ubicación & Potencia & Piso & SNR & Ancho 99\,\% & Capacidad \\
            & (dBFS) & (dBFS/Hz) & (dB) & (MHz) & (Mbps) \\
  \midrule
  Exterior    & $-$29.16 & $-$141.33 & 40.74 & 13.92 & 188.3 \\
  Laboratorio & $-$47.29 & $-$140.57 & 21.93 & 14.08 & 102.7 \\
  Pasillo     & $-$54.30 & $-$141.16 & 15.86 & 13.99 & 74.2 \\
  \bottomrule
\end{tabular}

\subsection{Portadoras detectadas}

\begin{tabular}{@{}lccc@{}}
  \toprule
  Ubicación & Portadora & Frecuencia (MHz) & Ancho (MHz) \\
  \midrule
  Exterior & 10 MHz & 878.54 & 8.88 \\
           & 5 MHz  & 885.38 & 4.45 \\
  Laboratorio & 10 MHz & 878.08 & 8.88 \\
              & 5 MHz  & 885.24 & 4.39 \\
  Pasillo & 10 MHz & 874.42 a 881.00 & 4 fragmentos \\
          & 5 MHz  & 883.60 a 887.26 & 3 fragmentos \\
  \bottomrule
\end{tabular}

En el exterior y el laboratorio se detectaron dos regiones, una por portadora. En el pasillo
la portadora de 10 MHz se partió en cuatro fragmentos (874.42, 876.60, 878.77 y
881.00 MHz) y la de 5 MHz en tres (883.60, 885.39 y 887.26 MHz). Esa fragmentación es
consecuencia de los nulos de desvanecimiento del canal, no de que la emisión haya cambiado.

\section{Comparación entre ubicaciones}

\figura{figs/comparacion.png}{Potencia total y SNR en banda para las tres ubicaciones.}

La potencia cae 18.1 dB del exterior al laboratorio y 7.0 dB adicionales al pasillo. La SNR
sigue la misma tendencia: 40.7, 21.9 y 15.9 dB. El piso de ruido es prácticamente el mismo
en las tres mediciones ($-$140.6 a $-$141.3 dBFS/Hz), lo que confirma que la diferencia se
debe a la atenuación y a la dispersión del canal, no al receptor.

La mayor degradación del pasillo proviene de dos efectos combinados: la atenuación por las
paredes y la multitrayectoria, que reparte la energía en ecos con retardos de hasta
15 muestras ($\approx$1 µs) y crea desvanecimientos profundos cada 1--2 MHz.

\section{Relación con OFDM y MIMO}

\figura{figs/estructura_ofdm.png}{Detalle de la portadora de 10 MHz (meseta característica de OFDM, inclinada por la respuesta del canal) y forma de onda temporal.}

La meseta con bordes abruptos es la misma que se obtuvo en la práctica 5 al generar
una señal OFDM con la IFFT: la suma de cientos de subportadoras de 15 kHz produce una
densidad espectral casi constante dentro del canal. En el tiempo, la envolvente tiene el
aspecto gaussiano de una suma de muchas componentes independientes, con picos que obligan a
reservar margen en el amplificador (práctica 8).

La estructura fina de 15 kHz no se resuelve en la PSD promediada, porque el ruido de
medición y el promediado de Welch la suavizan; para verla habría que demodular la señal con
un receptor LTE. Lo que sí se observa es la huella espectral de OFDM y la ocupación de
9 MHz y 4.5 MHz, compatible con 600 y 300 subportadoras.

Respecto de MIMO, la capacidad medida es la de un enlace SISO equivalente. Si la estación
base usara 2x2 o 4x4, la capacidad podría multiplicarse por 2 o 4 en las mismas
condiciones de SNR, como se mostró en la práctica 9. El límite práctico lo ponen la
correlación entre antenas, el scattering del entorno y la calidad del canal.

\section{Cálculo de capacidad}

\begin{tabular}{@{}lcccc@{}}
  \toprule
  Ubicación & SNR (dB) & SNR lineal & Eficiencia (bps/Hz) & Capacidad (Mbps) \\
  \midrule
  Exterior    & 40.74 & $1.19\times10^{4}$ & 13.53 & 188.3 \\
  Laboratorio & 21.93 & 156.1              &  7.29 & 102.7 \\
  Pasillo     & 15.86 & 38.6               &  5.31 & 74.2 \\
  \bottomrule
\end{tabular}

La eficiencia espectral cae de 13.53 a 5.31 bps/Hz, es decir, un factor 2.5 entre el
exterior y el pasillo. La capacidad sigue la misma proporción porque el ancho de banda
ocupado es prácticamente el mismo en las tres ubicaciones (13.9 a 14.1 MHz).

\section{Conclusiones}

Se caracterizó una red LTE de 850 MHz en tres escenarios. La potencia recibida varió de
$-$29.16 dBFS en el exterior a $-$54.30 dBFS en el pasillo, con SNR de 40.7 a 15.9 dB y
capacidades estimadas de 188.3 a 74.2 Mbps. El piso de ruido se mantuvo constante, lo que
atribuye las diferencias a la propagación y a la multitrayectoria. En el pasillo los ecos
fragmentaron el espectro de las portadoras en varios lóbulos, evidencia experimental de la
dispersión del canal. La forma plana del espectro y la ocupación de 9 y 4.5 MHz confirman
la naturaleza OFDM de LTE.

\section{Referencias}

\begin{enumerate}
  \item 3GPP TS 36.101, \emph{User Equipment (UE) radio transmission and reception},
        Release 15, 2018.
  \item T. S. Rappaport, \emph{Wireless Communications: Principles and Practice},
        2a ed., Prentice Hall, 2002.
  \item A. Goldsmith, \emph{Wireless Communications}, Cambridge University Press, 2005.
  \item V. I. Rodríguez Abdalá, \emph{Manual de Prácticas: Tecnologías Celulares},
        UAZ, 2026.
\end{enumerate}

\end{document}
